\section{Introduction}

Integer compositions are ways of writing an integer $n$ as a sum of positive integers where order matters. When zero is allowed to be a part of the sequence, they are called weak compositions.

Compositions have been studied for centuries and much of what we know today is around counting them with specific structures in place. The sheer volume of results in The On-Line Encyclopedia of Integer Sequences (OEIS) for compositions speaks to the depth to which these structures have been studied. Naturally, with the rapid advances in computing technology, generation of combinatorial objects, including compositions, has become more prevalent (Cite: A History of Algorithms). Knuth himself speaks to the explosion of algorithms, particularly combinatorial algorithms, with the rise of computers (Cite: Art of Programming Volume 4A). Beyond simply generating results, the ability to generate results in lexicographic order is particularly desirable as doing so allows us to create an isomorphism with the natural numbers, $\mathbb{N}$. This gives way to ranking and unranking algorithms, the latter providing a basis for parallel algorithms along with sampling.

At present, simply generating standard compositions in an unrestricted manner has been thoroughly covered. Because of the connection between compositions and binomial coefficients, illustrated beautifully by the balls-and-bars model (Cite: Analytic Combinatorics), implementations abound for generating compositions. Many of these start with the number to be decomposed, $n$, and produce successive results in reverse lexicographic order, with a general movement from top left to bottom right. From here, we can make small adjustments in order to produce compositions in lexicographic order. Instead of starting at the far left, place $n$ at the far right, and modify indexing to reflect a right to left movement. We can further alter these algorithms with minimal effort, for example by requiring the output to be of a specific length.

What if we restrict the output to distinct parts? To our knowledge, existing methods do not address this case. What made the algorithms for standard compositions so approachable was the fact that with each iteration, one part is incremented, while some other part is decremented. When the algorithm reaches the stage of incrementing, there is no required knowledge of the other parts; no need for global bookkeeping to avoid accidentally moving to a degenerate state. Adding the restriction of distinct parts changes all of this.

We present here an efficient algorithm for generating compositions with distinct parts in lexicographic order. The algorithm has time and space complexity of $\bigO(n)$. Following this, we also present ranking and unranking algorithms.
